%/*% pdflatex -file-line-error -interaction nonstopmode %


\documentclass{article}
\usepackage{graphicx}
\usepackage{amsmath}
\usepackage{amssymb}
\usepackage{amsfonts}
\usepackage{hyperref}
\usepackage{exercises}
\usepackage{amsthm}
\usepackage{yhmath}
\usepackage{placeins}
\newcommand*{\oldneg}{\mathord{\sim}}
\let\oldemptyset\emptyset

\theoremstyle{definition}
\newtheorem{definition}{Definition}

\setlength{\parindent}{.5in}
\setlength{\parskip}{12pt}
\title{Chapter 2 Answers}
\begin{document}
\maketitle

\section*{Lemma}
We must show that, from $F_1$ and $F_2$, the following two conditions hold:\newline
(1): We need to show that $\oldneg\exists x\phi(x)$ is true iff $\oldneg\phi(a)$ is true for every parameter $a$. Suppose that $\oldneg\exists x\phi(x)$ 
is true (under some valuation $v$). Thus, by $F_2$, the sentence $\phi(a)$ is not true under $v$ for any parameter $a$. Because $\phi(a)$ is not 
true under $v$ for any parameter $a$, it follows that $\oldneg\phi(a)$ is true under $v$ for every parameter $a$ (by $B_1$). On the other hand, 
suppose that $\oldneg\phi(a)$ is true under $v$ for every parameter $a$. Thus, by $B_1$, $\phi(a)$ is false under $v$ for every parameter $a$ (by $B_1$). Thus, by $F_2$, $\exists x\phi(x)$ is not true under $v$, meaning, by $B_1$, $\oldneg\exists x\phi(x)$ is true under $v$.\newline
(2): We need to show that $\oldneg\forall x\phi(x)$ is true iff $\oldneg\phi(a)$ is true for at least one parameter $a$. On the one hand, suppose 
that $\oldneg\forall x\phi(x)$ is true under some valuation $v$. Thus $\forall x\phi(x)$ is not true under $v$ (by $B_1$), meaning, by $F_1$, 
$\phi(a)$ is not true for every parameter $a$. Thus, for the parameter $e$ for which $\phi(e)$ is not true, $\oldneg\phi(e)$ is true by $B_1$. On the 
other hand, suppose that $\oldneg\phi(a)$ is true for at least one parameter $a$, call it $e$. Thus, by $B_1$, $\phi(e)$ is false, meaning, by $F_1$, 
$\forall x\phi(x)$ is not true, and therefore, by $B_1$, $\oldneg\forall x\phi(x)$ is true.  

\section*{1}
We suppose that $v$ is a Boolean valuation such that the following two conditions hold (for all sentences $\gamma$ and $\delta$):\newline
(1) if $\gamma$ is true under $v$, then for every parameter $a$, the sentence $\gamma(a)$ is true under $v$, and\newline
(2) if $\delta$ is true under $v$, then $\delta(a)$ is true under $v$ for at least one parameter $a$.\newline
We must prove that $v$ is a first-order valuation. Thus, we must derive $F_1$ and $F_2$ in the definition of a FOL-valuation.\newline
$F_1$: On the one hand, if $\forall x\phi(x)$ is true, we know that, for every parameter $a$, the sentence $\phi(a)$ is true (by (1)). On the other hand,
if for every parameter $a$, the sentence $\phi(a)$ is true, we know that $\oldneg\exists x\oldneg\phi(x)$ is true under $v$ (by (2) and $B_1$), meaning
$\oldneg\oldneg\phi(a)$ is true under $v$ for every parameter $a$, which means $\phi(a)$ is true for every parameter $a$ (by $B_1$).\newline
$F_2$: On the one hand, if $\exists x\phi(x)$ is true, we know that, by (2), $\phi(a)$ is true for at least one parameter $a$. On the other hand, 
if $\phi(a)$ is true for at least one parameter $a$, we know (by $B_1$) that $\oldneg\phi(a)$ is false for at least one parameter $a$, meaning, by 
(1), $\oldneg\exists x\phi(x)$ is not true, meaning  

\end{document}
